% Ziang Li's CV -- standalone pdfLaTeX source.
% Based on your edited CV.zip; all entered dates, grades, and roles are retained.
% Master thesis and supervisors are confirmed. Compile with pdfLaTeX.
\documentclass[11pt,a4paper]{article}
\usepackage[margin=19mm,top=17mm,bottom=19mm]{geometry}
\usepackage[T1]{fontenc}
\usepackage[utf8]{inputenc}
\usepackage{lmodern}
\usepackage{amssymb}
\usepackage{xcolor}
\usepackage{tabularx}
\usepackage{array}
\usepackage{titlesec}
\usepackage{fancyhdr}
\usepackage{needspace}
\usepackage{hyperref}

% ---------- Personal information ----------
\newcommand{\CVName}{Ziang Li}
\newcommand{\CVEmail}{Ziang.Li@stud.uni-regensburg.de}
\newcommand{\CVAffiliation}{University of Regensburg}
\newcommand{\MasterDates}{\daterange{Oct 2025}{Present}}
\newcommand{\BachelorDates}{\daterange{Sep 2021}{Jun 2025}}
\newcommand{\CVUpdated}{October 2026}

% ---------- Master thesis and supervisors ----------
\newcommand{\MasterThesis}{An Isomorphism Between the Blow-Up of the Deligne–Mumford Moduli Space and the Decorated Deligne–Mumford Moduli Space}
\newcommand{\MasterSupervisors}{\cvlink{https://friedl.app.uni-regensburg.de/}{Prof. Stefan Friedl} and \cvlink{https://ammann.app.uni-regensburg.de/}{Prof. Bernd Ammann}}

% ---------- Light-blue design ----------
\definecolor{SkyBlue}{HTML}{8DCCED}
\definecolor{PaleBlue}{HTML}{EDF7FD}
\definecolor{InkBlue}{HTML}{2D607D}
\definecolor{BodyInk}{HTML}{263746}
\definecolor{MutedInk}{HTML}{576A78}
\color{BodyInk}
\hypersetup{colorlinks=true,urlcolor=InkBlue,linkcolor=InkBlue,
  pdftitle={Curriculum vitae - Ziang Li},pdfauthor={Ziang Li}}
\urlstyle{same}
\setlength{\parindent}{0pt}
\setlength{\parskip}{1pt}
\renewcommand{\arraystretch}{1.02}
\linespread{0.98}
\titleformat{\section}{\large\sffamily\bfseries\color{InkBlue}}{}{0pt}{}
  [\vspace{-3pt}{\color{SkyBlue}\titlerule[0.7pt]}]
\titlespacing*{\section}{0pt}{9pt}{5pt}
\titleformat{\subsection}{\normalsize\sffamily\bfseries\color{InkBlue}}{}{0pt}{}
\titlespacing*{\subsection}{0pt}{5pt}{3pt}
\pagestyle{fancy}
\fancyhf{}
\renewcommand{\headrulewidth}{0pt}
\renewcommand{\footrulewidth}{0.4pt}
\fancyfoot[L]{\footnotesize\sffamily\color{MutedInk}\CVName\quad /\quad Curriculum vitae}
\fancyfoot[R]{\footnotesize\sffamily\color{MutedInk}\thepage}

% One date-range style; it wraps naturally only when the date column is narrow.
\newcommand{\daterange}[2]{\mbox{#1}\,--\allowbreak\ \mbox{#2}}
\newcommand{\cvlink}[2]{\href{#1}{#2}}

% Top-aligned entries: dates are raised slightly to match the title visually.
\newcommand{\cventry}[4]{%
  \Needspace{4\baselineskip}%
  \noindent\raisebox{0.5pt}{\begin{minipage}[t]{27mm}
    \vspace{0pt}\small\sffamily\color{MutedInk}\strut#1
  \end{minipage}}\hspace{4mm}%
  \begin{minipage}[t]{\dimexpr\linewidth-31mm\relax}
    \vspace{0pt}\small\strut\textbf{#2}%
    \if\relax\detokenize{#3}\relax\else\par#3\fi
    \if\relax\detokenize{#4}\relax\else\par{\color{MutedInk}#4}\fi
  \end{minipage}\par\vspace{0.5pt}%
}

% Less prominent, single-line reading-seminar entries; dates use a wider column.
\newcommand{\seminarentry}[3]{%
  \Needspace{2\baselineskip}%
  \noindent\raisebox{0.5pt}{\begin{minipage}[t]{41mm}
    \vspace{0pt}\footnotesize\sffamily\color{MutedInk}\strut#1
  \end{minipage}}\hspace{4mm}%
  \begin{minipage}[t]{\dimexpr\linewidth-45mm\relax}
    \vspace{0pt}\footnotesize\strut#2\enspace\textcolor{MutedInk}{(#3)}
  \end{minipage}\par\vspace{0pt}%
}

% Two coursework columns that wrap within the entry's width.
\newcommand{\coursework}[1]{%
  \begingroup
  \setlength{\tabcolsep}{0pt}%
  \begin{tabular}{@{}>{\raggedright\arraybackslash}p{0.57\linewidth}@{\hspace{0.04\linewidth}}>{\raggedright\arraybackslash}p{0.39\linewidth}@{}}
    #1
  \end{tabular}%
  \endgroup
}

\begin{document}

\begingroup
\setlength{\fboxsep}{10pt}
\colorbox{PaleBlue}{\begin{minipage}{\dimexpr\linewidth-20pt\relax}
  {\fontsize{27}{30}\selectfont\sffamily\bfseries\color{InkBlue}\CVName}\par
  \vspace{2pt}
  {\sffamily Master's student in mathematics}\par
  {\small\CVAffiliation\quad |\quad Regensburg, Germany}\par
  \vspace{2pt}
  {\small\href{mailto:\CVEmail}{\CVEmail}}
\end{minipage}}
\endgroup

\section{Education}
\cventry{\MasterDates}{University of Regensburg}
  {Master's studies in mathematics; Regensburg, Germany}
  {\textbf{Master thesis in progress:} \MasterThesis.\par
   \textbf{Supervisor(s):} \MasterSupervisors\par
   \textbf{Selected coursework:}\par
   \textit{Grades in Germany are on a scale from 1.0 to 5.0; 5.0 is a failing grade.}\par
   \coursework{
     Twisted homology and Reidemeister torsion (1.3) & Algebraic Topology I (1.3)\\
     Simple homotopy theory and manifold topology (1.0) & Algebraic Topology II (1.0)\\
     Seminar on Floer homology, Part II (1.0) & Characteristic classes (1.3)
   }}
\cventry{\BachelorDates}{Xi'an Jiaotong University}
  {Bachelor's degree in mathematics; Xi'an, China}
  {\textbf{Thesis:} Quasi-isometry Groups of $\mathbb{R}^n$.\par
   \textbf{Supervisor:} Qiang Zhang.\par
   \textbf{Selected coursework:}\par
   \coursework{
     Riemannian Geometry (100) & Functional Analysis (90)\\
     Graph Theory and Combinatorics (92) & Algebra Seminar (92)\\
     Mathematical Analysis Seminar (99) & Differentiable Manifolds (90)
   }}

\section{Research interests}
{\small Geometric topology, especially knot theory and the topology of 3- and 4-manifolds.}

\section{Academic experience}
\subsection{Seminar talks}
% The report links below point to the same PDFs used on the homepage.
% Keep the compiled CV PDF in files/, beside these four report PDFs.
% For a CV shared separately, replace these relative paths with public URLs.
\cventry{Jul 2026}{Proof of the $s$-cobordism theorem}
  {Simple homotopy theory and manifold topology seminar, University of Regensburg}
  {Organizer: Dr. Marco Volpe.\quad\cvlink{rep-scob.pdf}{Report}}
\cventry{Jul 2026}{A surgery formula for Seiberg--Witten invariants}
  {Seminar on spin geometry, University of Regensburg}
  {Organizer: Prof. Bernd Ammann.\quad\cvlink{rep-Surgery.pdf}{Report}}
\cventry{Dec 2025}{Proof of the Arnold conjecture}
  {Seminar on Floer homology, Part II, University of Regensburg}
  {Organizer: Prof. Bernd Ammann.\quad\cvlink{rep-Floer.pdf}{Report}}
\cventry{Nov 2025}{An introduction to Seiberg--Witten theory and its applications to 4-manifolds}
  {Surfaces in 4-manifolds seminar, University of Regensburg}
  {Organizer: Prof. Stefan Friedl.\quad\cvlink{rep-SW.pdf}{Report}}

\section{Research notes}
% Notes are listed in the same order as on the homepage.
\cventry{Jul 2026}{Direct proof of equivalence of two definitions of Alexander polynomials}
  {Supervised by: \cvlink{https://friedl.app.uni-regensburg.de/}{Prof. Stefan Friedl}.}
  {I gave a direct proof that the definitions of the Alexander polynomial via the Burau representation and Fox calculus are equivalent.}
\cventry{Sep 2026}{Master thesis in progress: An Isomorphism Between the Blow-Up of the Deligne–Mumford Moduli Space and the Decorated Deligne–Mumford Moduli Space.}
  {Supervised by: \cvlink{https://ammann.app.uni-regensburg.de/}{Prof. Bernd Ammann}.}
  {Tba}

\section{Conferences, summer schools, and workshops}
% Program links are copied from the homepage; these entries record participation.
\cventry{\daterange{8 Jun 2026}{12 Jun 2026}}{Masterclass: Seiberg-Witten Invariants}
  {University of Copenhagen; Copenhagen, Denmark.\quad\cvlink{https://www.math.ku.dk/english/calendar/events/masterclass-seiberg-witten-theory/}{Program}}{}
\cventry{\daterange{18 May 2026}{22 May 2026}}{Spring school "Khovanov and Floer homotopy"}
  {Nantes University; Nantes, France.\quad\cvlink{https://springschoolhomotopy-fabiogironella.apps.math.cnrs.fr/}{Program}}{}
\cventry{\daterange{27 Dec 2025}{28 Dec 2025}}{ALGANT Alumni in China}
  {University of Science and Technology of China; Hefei, China.\quad\cvlink{https://faculty.ustc.edu.cn/yqliang/files/ALGANTinChina.htm}{Program}}{}
\cventry{\daterange{8 Dec 2025}{12 Dec 2025}}{Workshop on surfaces in 4-manifolds}
  {University of Regensburg; Regensburg, Germany.\quad\cvlink{https://sites.google.com/view/surfacesin4-manifolds/home}{Program}}{}
\cventry{\daterange{1 Jul 2025}{11 Jul 2025}}{USTC's 2025 Summer School in Geometry and Analysis}
  {University of Science and Technology of China; Hefei, China.\quad\cvlink{https://igp.ustc.edu.cn/2025/0403/c28839a679308/page.htm}{Program}}{}
\cventry{\daterange{22 Aug 2024}{28 Aug 2024}}{SUSTech’s 2024 Number Theory and Arithmetic Geometry Conference}
  {Southern University of Science and Technology; Shenzhen, China.\quad\cvlink{https://math.sustech.edu.cn/conference/13077.html}{Program}}{}
\cventry{\daterange{29 Jul 2024}{9 Aug 2024}}{BICMR’s 2024 Summer School in Differential Geometry}
  {Beijing International Center for Mathematical Research, Peking University; Beijing, China.\quad\cvlink{https://bicmr.pku.edu.cn/cn/content/show/17-3335.html}{Program}}{}
\cventry{\daterange{1 Jul 2024}{9 Jul 2024}}{YMSC’s 2024 Summer School on Frontiers in Topology}
  {Yau Mathematical Sciences Center, Tsinghua University; Beijing, China.\quad\cvlink{https://ymsc.tsinghua.edu.cn/info/1060/3688.htm}{Program}}{}

\section{Student-led reading seminars}
\seminarentry{\daterange{Sep 2024}{Jun 2025}}{Symplectic geometry}{Organizer and main speaker}
\seminarentry{\daterange{Mar 2024}{Jun 2024}}{Characteristic classes}{Organizer and main speaker}
\seminarentry{\daterange{Dec 2023}{Mar 2024}}{Morse theory}{Organizer and main speaker}
\seminarentry{\daterange{Sep 2023}{Jun 2024}}{Category theory and homological algebra}{Participant}
\seminarentry{\daterange{Mar 2022}{Jan 2023}}{Algebraic topology}{Participant and speaker}
\seminarentry{\daterange{Feb 2022}{Jun 2022}}{Riemannian geometry}{Organizer and main speaker}
\seminarentry{\daterange{Sep 2021}{Jan 2022}}{Differential manifolds and de Rham cohomology}{Organizer and main speaker}

\vfill
{\footnotesize\sffamily\color{MutedInk}Last updated: \CVUpdated.}
\end{document}
